\subsubsection{Existence of Weak Solutions}\label{subsec:dirichlet-weak-solution-existence}

We are now in a position to prove our basic existence result for weak solutions.

\begin{Theorem}
\label{thm:weak-dirichlet-positive-shift-existence}
Let $L(\bm{x},D)$ be a linear partial differential operator in divergence form of order $2k$, satisfying the hypotheses of Theorem~\ref{thm:garding-elliptic-energy-inequality} (Gårding\textquotesingle s inequality). Then there exists $\lambda_G\geq0$ such that for any $\widetilde\lambda\geq\lambda_G$, and any $f\in H^{-k}(\Omega)$, the Dirichlet problem for the operator
\begin{equation}\label{eq:weak-dirichlet-positive-shifted-operator}
\widetilde L(\bm{x},D):=L(\bm{x},D)+\widetilde\lambda
\end{equation}
has a unique weak solution $u\in H_0^k(\Omega)$. Furthermore, this solution satisfies
\begin{equation}\label{eq:weak-dirichlet-solution-negative-norm-bound}
\|u\|_{k,2}\leq C\|f\|_{-k,2}.
\end{equation}

\end{Theorem}

\begin{Proof}
Theorem~\ref{thm:garding-elliptic-energy-inequality} guarantees the existence of $\lambda_G\geq0$ such that \eqref{eq:garding-elliptic-energy-inequality} holds. Let $\widetilde\lambda\geq\lambda_G$. Note that
\begin{equation}\label{eq:weak-dirichlet-shifted-bilinear-form}
\widetilde B[u,v]:=B[u,v]+\widetilde\lambda(u,v)_{L^2(\Omega)}
\end{equation}
is the bilinear form associated with the operator $\widetilde L$ defined in \eqref{eq:weak-dirichlet-positive-shifted-operator}. We now show that $\widetilde B$ satisfies the hypotheses of the Lax-Milgram lemma (Theorem~\ref{thm:lax-milgram}).

Let $H=H_0^k(\Omega)$, and let $u,v\in H$. Then\[\begin{aligned}
|\widetilde B[v,u]|&\leq|B[v,u]|+|\widetilde\lambda|(u,v)\\
&\leq\sum_{0\leq|\sigma|,|\gamma|\leq k}\int_\Omega|a_{\sigma\gamma}(\bm{x})|\,|D^\gamma u|\,|D^\sigma v|\,\mathrm{d}\bm{x}+|\widetilde\lambda|(u,v)\\
&\leq\max_{|\sigma|,|\gamma|\leq k}\|a_{\sigma\gamma}\|_{L^\infty(\Omega)}
\sum_{0\leq|\sigma|,|\gamma|\leq k}\int_\Omega|D^\gamma u|\,|D^\sigma v|\,\mathrm{d}\bm{x}\\
&\quad+|\widetilde\lambda|(u,v)\\
&\leq C\|v\|_H\|u\|_H.
\end{aligned}\]Thus, $\widetilde B$ satisfies \eqref{eq:lax-milgram-bilinear-boundedness}.

Now by Gårding\textquotesingle s inequality \eqref{eq:garding-elliptic-energy-inequality} we have
\begin{equation}\label{eq:weak-dirichlet-shifted-form-coercivity}
\widetilde B[u,u]=\widetilde\lambda\|u\|_2+B[u,u]\geq c_3\|u\|_H^2.
\end{equation}
Thus, $\widetilde B$ satisfies \eqref{eq:lax-milgram-bilinear-coercivity}.

Thus, Lax-Milgram guarantees that for every $f\in H^{-k}=H^*$ there is a unique weak solution $u\in H$ of the Dirichlet problem, and that the solution satisfies the estimate \eqref{eq:weak-dirichlet-solution-negative-norm-bound}.
\end{Proof}
